%%%%%%%%%%%%%%%%%%%%%%%%%%%%%%%%%%%%%%%%%%%%%%%%%%%%
% Artifact Appendix Template for Usenix Security'25 AE
%%%%%%%%%%%%%%%%%%%%%%%%%%%%%%%%%%%%%%%%%%%%%%%%%%%%

% \appendix

\section{Artifact Appendix}
This artifact appendix provides a roadmap for evaluating the artifact
associated with our paper, ``When Authorization Loses Its Meaning: Breaking and 
Fixing Third-Party Online Payments''. 

%%%%%%%%%%%%%%%%%%%%%%%%%%%%%%%%%%%%%%%%%%%%%%%%%%%%%%%%%%%%%%%%%%%%%
\subsection{Abstract}
This artifact provides the formal models and verification
framework used to analyze the security of third-party online payment protocols.

%%%%%%%%%%%%%%%%%%%%%%%%%%%%%%%%%%%%%%%%%%%%%%%%%%%%%%%%%%%%%%%%%%%%%
\subsection{Description \& Requirements}
This section describes the steps and requirements needed to set up the artifact and reproduce the experiments.

\subsubsection{Security, privacy, and ethical concerns}
None. The formal models in this artifact are derived exclusively from publicly available developer documentation of the studied payment platforms.
No reverse engineering, no inspection of proprietary or internal documents, and no testing against live accounts or production systems was performed.

\subsubsection{How to access}
The artifact is archived on Zenodo and can be accessed through the following DOI: \url{https://doi.org/10.5281/zenodo.20303820}.

\subsubsection{Hardware dependencies}
An x86-64 machine with at least 8 CPU cores, 32~GB of RAM, and 10~GB of free disk space.

\subsubsection{Software dependencies}
Ubuntu 24.04 with Docker installed. 

\subsubsection{Benchmarks}
None


%%%%%%%%%%%%%%%%%%%%%%%%%%%%%%%%%%%%%%%%%%%%%%%%%%%%%%%%%%%%%%%%%%%%%
\subsection{Set-up}

This section describes the installation and basic test steps required to prepare
the environment for evaluating the artifact.

\subsubsection{Installation}
\begin{enumerate}
    \item Pull our pre-built Docker image.
\begin{verbatim}
$ docker pull ghcr.io/luojiazhishu\
/tamarin-docker/toolbox:latest
\end{verbatim}
    \item Download and extract the artifact archive from Zenodo (\url{https://doi.org/10.5281/zenodo.20303820}). 
    After extraction, the directory structure is as follows:
\begin{Verbatim}[baselinestretch=0.85]
.
|-- models
|   |-- base
|   |-- countermeasures
|   |-- secure_ms_channel
|   `-- secure_ps_channel
|-- results
|   |-- base
|   |-- countermeasures
|   |-- secure_ms_channel
|   `-- secure_ps_channel
|-- scripts
|   |-- fix.py
|   |-- parse.py
|   `-- run.py
|-- Makefile
`-- README.md
\end{Verbatim}
\end{enumerate}

The top-level entries are organized as follows:
\begin{itemize}[itemsep=1pt, topsep=2pt, parsep=0pt, partopsep=0pt]
    \item \texttt{./models/}: Tamarin formal models. The subdirectories \texttt{secure\_ms\_channel/}, \texttt{secure\_ps\_channel/}, and \\
    \texttt{base/} correspond to three channel security assumptions,
    each containing six payment-protocol models (two payment providers $\times$ three payment scenarios) --- 18 protocol models in total. The formal models of our proposed countermeasures are under \texttt{countermeasures/}.
    \item \texttt{./scripts/}: Helper scripts for running Tamarin and aggregating its output into JSON/CSV summaries.
    \item \texttt{./results/}: Verification results produced by Tamarin, mirroring the subdirectory layout of \texttt{./models/}.
    \item \texttt{./Makefile}: Entry point that orchestrates the full evaluation pipeline.
\end{itemize}

\subsubsection{Basic Test}
To verify the environment, run:

\begin{verbatim}
$ docker run --rm ghcr.io/luojiazhishu\
/tamarin-docker/toolbox:latest \
  tamarin-prover test
\end{verbatim}

If the installation is successful, the command output contains the following test summary:
\begin{verbatim}
*** TEST SUMMARY ***
All tests successful.
The tamarin-prover should work as intended.

           :-) happy proving (-:
\end{verbatim}


%%%%%%%%%%%%%%%%%%%%%%%%%%%%%%%%%%%%%%%%%%%%%%%%%%%%%%%%%%%%%%%%%%%%%
\subsection{Evaluation workflow}
The evaluation workflow involves formal model verification using the Tamarin Prover.

\subsubsection{Major Claims}

\begin{compactdesc}

    \item[(C1):] The verification results of six third-party online payment protocols under three channel assumptions reproduce Table~1 in the paper. Supported by experiment \textbf{(E1)}.

    \item[(C2):] The verification results of our countermeasures reproduce Table~2 in the paper. Supported by experiment \textbf{(E1)}.


\end{compactdesc}

\subsubsection{Experiments}

% use paralist for more compact list format: for more details check here:
% https://texfaq.org/FAQ-complist
\begin{compactdesc}

    \item[(E1):] \textbf{Verification} [ 3 human-minutes + 1 compute-hour ]: Use the Tamarin Prover to check the security lemmas of every formal model in \texttt{./models/}. The resulting \textit{verified}/\textit{falsified} outcomes should match Table~1 and Table~2 of the paper.

    \begin{asparadesc}

        \item[How to:] Run the verification in three steps.

       \item[Preparation:]
        With the installation in Section~A.3 completed, launch the Tamarin toolbox container from the \textbf{artifact root directory}:
\begin{verbatim}
$ docker run -d \
  --name tamarin-toolbox \
  -p 3001:3001 \
  -v "$(pwd):/root" \
  -w /root \
  ghcr.io/luojiazhishu\
/tamarin-docker/toolbox:latest
\end{verbatim}

        \item[Execution:]
        Enter the container, clear any pre-existing results, and run the verification pipeline:
\begin{verbatim}
$ docker exec -it -w /root \
  tamarin-toolbox bash
$ make clean
$ make all JOBS=1
\end{verbatim}
        On the reference hardware (8 cores, 32~GB RAM), \texttt{JOBS=1} completes the full verification in about one hour. On a more powerful machine (e.g., 32 cores and 64~GB RAM), increase parallelism with \texttt{make all JOBS=8} to reduce the total time to roughly 30 minutes.

        \item[Results:]
        The top-level file \texttt{./results/results.csv} aggregates the verification outcome of every lemma across all models; its rows reproduce Table~1 and Table~2 of the paper. Each subdirectory of \texttt{./results/} corresponds to one channel security assumption (or to the countermeasures) and contains one \texttt{.spthy} proof file per protocol plus a \texttt{result.json} summarising that subdirectory. To inspect a proof graph, launch Tamarin's web interface inside the container; for example, to view the proofs under \texttt{./results/base/}, run:
\begin{verbatim}
$ tamarin-prover interactive \
  --interface=0.0.0.0 \
  --derivcheck-timeout=0 \
  ./results/base
\end{verbatim}
        Then open \texttt{http://127.0.0.1:3001} in a browser on the host machine.

    \end{asparadesc}

\end{compactdesc}

\noindent The Tamarin lemma names that appear in \texttt{results.csv} and in each \texttt{result.json} correspond to the security properties defined in the paper as follows:
\begin{itemize}[itemsep=0pt, topsep=1pt, parsep=0pt, leftmargin=*]
    \item \texttt{PS\_MS\_Order\_Injective\_Agreement: IA-O1}.
    \item \texttt{U\_MS\_Order\_Injective\_Agreement: IA-O2}.
    \item \texttt{MS\_PS\_Response\_Injective\_Agreement: IA-R1}.
    \item \texttt{U\_PS\_Response\_Injective\_Agreement: IA-R2}.
    \item \texttt{No\_Payment\_Without\_Password\_Input: IA-UA}.
\end{itemize}

%%%%%%%%%%%%%%%%%%%%%%%%%%%%%%%%%%%%%%%%%%%%%%%%%%%%%%%%%%%%%%%%%%%%%

\subsection{Version}
%%%%%%%%%%%%%%%%%%%%
% Obligatory.
% Do not change/remove.
%%%%%%%%%%%%%%%%%%%%
Based on the LaTeX template for Artifact Evaluation V20231005. Submission,
reviewing and badging methodology followed for the evaluation of this artifact
can be found at \url{https://secartifacts.github.io/usenixsec2026/}.